\emph{Setting.} The linearisation is that of Theorem \claimref{synthesis.horizon-channels-of-the-linearised-gravitational-field-are-two-dimensional}: the static coframe of Fact \claimref{pairing.imported-schwarzschild-connection-curvature}, $g_{\mu} \mapsto g_{\mu} + \epsilon k_{\mu}$, $\Omega(e_{\mu}) \mapsto \Omega(e_{\mu}) + \epsilon\Omega_{1}(e_{\mu})$, $\epsilon^{2} = 0$; $\Omega_{1}$ is the unique solution of the first-order torsion-free condition $\Omega_{1}(e_{\mu}) \cdot g_{\nu} - \Omega_{1}(e_{\nu}) \cdot g_{\mu} = -(D_{\mu}k_{\nu} - D_{\nu}k_{\mu})$ (\claimref{obj.synthesis.gtg-flat-spacetime-gauge-fields}), and $\mathcal{R}_{1}(\gamma_{b})$ is the first-order part of the Ricci vector of the curvature operator in the store's sign. The stated $k_{\mu}$ give the first-order metric $-fH_{0}P\,dt^{2} - 2H_{1}P\,dt\,dr - f^{-1}H_{2}P\,dr^{2} - r^{2}KP\,(d\vartheta^{2} + \sin^{2}\vartheta\,d\varphi^{2})$.

\emph{Why a basis of jets decides (1) and (2).} $\Omega_{1}$ at a point is a linear function of $k$ and its first derivatives there (the torsion system is algebraic in $\Omega_{1}$, with the invertible matrix of the background torsion map), and $\mathcal{R}_{1}$ is a linear function of $\Omega_{1}$, its first derivatives and $k$; so $\mathcal{R}_{1}(\gamma_{b})$ at a point is a linear function of the second-order jet of $k$ at that point. For the even-parity field each component of $k_{\mu}$ is a background function times a product $H(t,r)P(\vartheta)$, whose second-order jet at the point is a bilinear function of the second-order jets of $H$ and of $P$ there. The displayed right-hand sides are bilinear in the same jets. So the two sides agree for all smooth $H_{0}, H_{1}, H_{2}, K, P$ at the point as soon as they agree on a basis: the six monomials $t^{p}(r - r_{0})^{q}$, $p + q \leq 2$, for each of the four functions, and the three monomials $1$, $\vartheta - \vartheta_{0}$, $\tfrac12(\vartheta - \vartheta_{0})^{2}$ for $P$. Nothing is assumed of $P$: the term $P''$ is independent data, so no separation constant and no degree count in $L$ or in $\cot\vartheta$ enters.

\emph{Why the comparison is an identity in the radius and the angle.} The check runs the linearisation with jets whose coefficients lie in $U = \mathbb{Q}[z, z^{-1}, (1 - z^{2})^{-1}, s, s^{-1}, c]/(c^{2} + s^{2} - 1)$, the base point being $r_{0} = 2/(1 - z^{2})$, $\sqrt{f(r_{0})} = z$, $\sin\vartheta_{0} = s$, $\cos\vartheta_{0} = c$ ($M = 1$). An element of $U$ is stored as a polynomial numerator, reduced by $c^{2} = 1 - s^{2}$ to degree at most one in $c$, over a monomial $z^{a}(1 - z^{2})^{b}s^{e}$; such a representation is zero only if the numerator is, since $z$, $1 - z^{2}$ and $s$ are not zero divisors. The implementation inverts only rational multiples of monomials in $z$, $1 - z^{2}$ and $s$ and raises an error otherwise; the computation completes, so every division made was by such a unit (the pivots of the torsion system are monomials in the frame components $z$, $1/z$, $r_{0}$, $r_{0}s$, as its homogeneity under independent rescalings of the four frame vectors requires). $U$ maps to the functions on $\{2M < r,\ \sin\vartheta \neq 0\}$ by evaluation, and the units are nonzero there, so an equality in $U$ is an equality of functions at every radius of the exterior and every such angle. All sixteen components are compared in each of the $72$ cases, the vanishing ones included ($72$ are nonzero). General $M$: every term of the closed forms is homogeneous under $(t, r, M) \mapsto (at, ar, aM)$ with $H_{0}, H_{1}, H_{2}, K$ of weight $0$ (weight $-2$ for $C_{00}, C_{01}, C_{11}, C_{22}$ and the $\Delta P$ terms, $-1$ for $C_{02}, C_{12}$, whose prefactor $1/(2r\sqrt f)$ has weight $-1$), and so is $\mathcal{R}_{1}$, the constant dilation of the chart being a displacement (clause (e) of the Definition) that carries the field of mass $M$ to a constant multiple of the field of mass $aM$; a sampled control at masses $\tfrac98$ and $\tfrac{25}{18}$ with generic jets is in the check. This proves (1).

\emph{(2).} The same computation with the odd-parity field of the cited theorem, $k_{\varphi} = S(f^{-1/2}h_{0}\gamma_{0} - f^{1/2}h_{1}\gamma_{1})$: its closed forms contain $L$ through $E_{t}$ and $E_{r}$, and they were derived for $S'' = \cot\vartheta\,S' - LS$. Both sides are linear in the second-order jet of $S$, hence, with $S''$ so expressed, of degree one in $L$ for fixed $(S, S')$; they agree in $U$ for $L = \tfrac72$ and $L = 6$ on the $12$ basis jets of $(h_{0}, h_{1})$ and the two basis data of $(S, S')$ ($48$ cases, $56$ nonzero components counting the symmetric partners), hence for every $L$, radius and angle.

\emph{(3).} With $\Delta P = -LP$ the components of (1) are $P E_{1}$, $\mp P E_{2}$, $\pm\tfrac{P'}{2r\sqrt f}E_{3}$, $\tfrac{P'}{2r\sqrt f}E_{4}$, $P E_{5}$, and $P E_{7} + \tfrac{1}{2r^{2}}E_{6}P''$, $P E_{7} + \tfrac{1}{2r^{2}}E_{6}\cot\vartheta\,P'$, where $E_{j}$ depend on $(t, r)$ only. If all $E_{j}$ vanish, so does $\mathcal{R}_{1}$. Conversely let $\mathcal{R}_{1} = 0$. $P$ is not identically zero, and neither is $P'$ since $L \neq 0$; so $E_{1} = E_{2} = E_{3} = E_{4} = E_{5} = 0$. The difference of the last two components is $\tfrac{1}{2r^{2}}E_{6}(P'' - \cot\vartheta\,P')$. If $P'' = \cot\vartheta\,P'$ identically, the equation of $P$ gives $2\cot\vartheta\,P' = -LP$, so $P = C\cos^{n}\vartheta$ with $n = L/2$ on an interval where $\cos\vartheta \neq 0$, and then $P'' + \cot\vartheta\,P' + LP = n(n-1)C\cos^{n-2}\vartheta\,\sin^{2}\vartheta$, which vanishes identically only for $n \in \{0, 1\}$, that is $L \in \{0, 2\}$. So for $L \notin \{0, 2\}$, $E_{6} = 0$ and then $E_{7} = 0$.

\emph{The differential algebra used for (4) and (5).} Let $A = \mathbb{Q}[r, \lambda, w][r^{-1}, (r - 2)^{-1}, (\lambda r + 3)^{-1}, w^{-1}]$, with $d/dr$ acting in the evident way, and consider linear forms with coefficients in $A$ in symbols $\partial_{t}^{p}\partial_{r}^{q}X$ for named functions $X$; $\partial_{r}$ and $\partial_{t}$ act on them by the Leibniz rule. A linear form is a linear differential expression in the functions $X$, and it is the zero form only if all its coefficients are zero in $A$ (numerators compared after bringing to a common denominator). The check evaluates the closed forms of (1) on such linear forms with $M = 1$, $f = (r - 2)/r$; every identity below is an equality of linear forms, hence an identity for arbitrary smooth functions at every $r > 2$, $\lambda > 0$, $w > 0$, where the inverted quantities do not vanish. Homogeneity as in (1), with $\Phi$ and $\Psi$ of weight $1$, $\omega$ of weight $-1$ (so $w = \omega^{2}$ of weight $-2$) and $V_{\mathrm{e}}$ of weight $-2$, gives general $M$: with these weights $a$ has weight $0$, $b$ weight $-1$, $\mathcal{Z}[\Phi]$ weight $-1$, $\mathrm{Rel}$ weight $1$, and each displayed identity is homogeneous.

\emph{(4).} Put $Z := \mathcal{Z}[\Phi]$, so that $\partial_{t}^{2}\Phi = f(f\Phi')' - V_{\mathrm{e}}\Phi + Z$, and substitute this (and its $t$- and $r$-derivatives) for every symbol $\partial_{t}^{p}\partial_{r}^{q}\Phi$ with $p \geq 2$ in the linear forms $E_{j}$ of the stated field. The substitution is an identity for every smooth $\Phi$, with $Z$ standing for $\mathcal{Z}[\Phi]$. The result, computed by the check: no symbol of $\Phi$ remains in any $E_{j}$; $E_{2}$, $E_{3}$, $E_{6}$ are zero (for $E_{3}$: $fH_{1} = \partial_{t}(fr\Phi' + fa\Phi)$ and $H_{2} + K$ are arranged so that $\partial_{r}(fH_{1}) = \partial_{t}(H_{2} + K)$; they contain no second time derivative, so the substitution is not used); $E_{4} = -(rZ' + aZ)$ and $E_{7} = -\tfrac12 Z' + bZ$ with the stated $a$, $b$; $E_{1}$ and $E_{5}$ are combinations of $Z$, $Z'$, $Z''$ with coefficients in $A$ and coefficient $-r/2$ of $Z''$. The Zerilli-Moncrief function of the field is computed directly, without substitution: $r\bigl(K + \tfrac{fr}{\lambda r + 3}(H_{2} - rK')\bigr) = (\lambda + 1)\Phi$ as linear forms (the terms in $\Phi'$ and $\Phi''$ cancel). With a potential differing by $f/r^{2}$ a term in $\Phi$ is left in $E_{4}$, so the test is sensitive to the potential. If $\mathcal{Z}[\Phi] = 0$ all $E_{j}$ vanish and by (3) the field is a vacuum solution; $\mathcal{Z}$ is linear, so $\Phi = \Psi/(\lambda + 1)$ solves the Zerilli equation when $\Psi$ does, and $\lambda + 1 = L/2 \neq 0$.

\emph{(5).} Let the field solve $E_{1} = \cdots = E_{7} = 0$ on the channel and write $H$ for $H_{0} = H_{2}$ ($E_{6}$). As linear forms, $E_{3} = f\,(H_{1}' - \rho_{1})$ and $E_{4} = f\,\bigl((H' - K') - (\rho_{H} - \rho_{K})\bigr)$, and, using $\partial_{t}^{2} = -\omega^{2}$, $\partial_{t}E_{2} = -\omega^{2}(K' - \rho_{K})$, where $\rho_{K}, \rho_{H}, \rho_{1}$ are the right-hand sides of the displayed first-order system (checked as equalities of linear forms). Since $f \neq 0$ and $\omega^{2} \neq 0$, the solution satisfies the system (only $\omega^{2} \neq 0$ is used, here and below; $\omega$ real is the oscillating case). Call \emph{reduction} the replacement of every symbol $\partial_{t}^{p}\partial_{r}^{q}X$, $X \in \{K, H, H_{1}\}$, with $q \geq 1$ by the corresponding derivative of the right-hand side, and of $\partial_{t}^{2}$ by $-\omega^{2}$, until only $K, H, H_{1}, \partial_{t}K, \partial_{t}H, \partial_{t}H_{1}$ remain; on the solution every expression equals its reduction. The check computes: $E_{2}, E_{3}, E_{4}, E_{6}$ reduce to zero; $E_{7}$ reduces to $\mathrm{Rel}/(2r^{3})$, a combination of $H$, $K$ and $\partial_{t}H_{1}$ with the stated coefficients, in which the coefficient of $H$, $-(\lambda r + 3)/(2r^{3})$, is a unit of $A$; $E_{1}$, $E_{5}$, $\partial_{r}\mathrm{Rel}$ and $\mathcal{Z}[\Psi]$ each reduce to $\mu\,\mathrm{Rel}$ with $\mu \in A$ (the quotient of the coefficients of $H$), a function regular on $r > 2$. So on the solution $\mathrm{Rel} = 0$ and $\mathcal{Z}[\Psi] = 0$: the Zerilli equation follows from the vacuum equations, identically. Finally $(\lambda + 1)K - K[\Psi]$, $(\lambda + 1)H - H[\Psi]$ and $(\lambda + 1)H_{1} - H_{1}[\Psi]$, where $K[\cdot], H[\cdot], H_{1}[\cdot]$ are the operators of (4), reduce to zero, without use of $\mathrm{Rel}$; so on the solution $(\lambda + 1)(K, H, H_{1})$ is the field of (4) built from $\Psi$, and $\Psi = 0$ forces the field to vanish. With (4): the map to $\Psi$ lands in the solutions of the Zerilli equation of the channel, is injective, and is onto with the stated inverse, because the field of (4) built from a channel solution $\Phi$ lies in the channel (its operators commute with $\partial_{t}$), solves the vacuum equations and has master function $(\lambda + 1)\Phi$. (That $E_{1}$, $E_{5}$ and $\partial_{r}\mathrm{Rel}$ are multiples of $\mathrm{Rel}$ is not needed for the bijection; it records that the first-order system with the one relation is the whole content of the seven equations on the channel: by the equalities above, the system gives $E_{3} = E_{4} = \partial_{t}E_{2} = 0$, hence $E_{2} = 0$ since $\omega \neq 0$, and with $\mathrm{Rel} = 0$ it gives $E_{1} = E_{5} = E_{7} = 0$.)

\emph{(6).} By (3) and (5), in the Regge-Wheeler gauge an even-parity vacuum solution on an oscillating channel with $L = \ell(\ell+1)$, $\ell \geq 2$, is the field of (4) built from a solution of $\mathcal{Z}[\Psi] = 0$, and conversely by (4). Clause (5) of the cited theorem gives $V_{\mathrm{e}}/f$ regular and positive at $r = 2M$ and $V_{\mathrm{e}} = O(e^{2\kappa r_{*}})$ as $r_{*} \to -\infty$, so $\Psi$ obeys the free two-dimensional equation there, as stated in that clause conditionally. The cited fact (C3), as stored, is the surjectivity alone: every solution of the Zerilli equation in an oscillating channel is the $\Psi$ of an even-parity vacuum solution; that is clause (4), which gives it without a condition on the time dependence. The forward implication and the injectivity, which the stored clause (5) also needs, are clause (5) here, and they are proved on channels with $\partial_{t}^{2} = -\omega^{2}$, $\omega^{2} \neq 0$, only: the stored clause (3) states the forward implication without that condition, at sample points, and in that generality it is not proved here. The reductions (C1) and (C2), from a general linearised field to the axisymmetric representative of one harmonic in the Regge-Wheeler gauge, are not touched by this proof; they enter (6) only through the phrase ``every even-parity channel'' and remain cited from memory.
